\section{Related Work}
\label{sec:related_work}

% ─────────────────────────────────────────────────────────────────────────────
\subsection{Continual Learning and the Stability-Plasticity Dilemma}

Modern machine learning systems are predominantly trained on fixed datasets
drawn from a single stationary distribution --- a condition that rarely holds
in deployment, where agents must adapt to non-stationary data streams while
retaining prior knowledge.
The fundamental tension between these two objectives is formalized as the
\emph{stability-plasticity dilemma} \citep{grossberg_how_1980}: an agent must
remain plastic enough to absorb novel distributions yet stable enough to resist
the erasure of established representations.


The continual learning literature addresses this across three canonical
scenarios \citep{de_lange_continual_2022}:
\emph{task-incremental learning}, where the task identity is available at
both training and test time; \emph{class-incremental learning}, where task
identity is withheld at inference; and \emph{domain-incremental learning},
where the output space is shared but the input distribution shifts
continuously.
This thesis focuses on the task-incremental setting, the scenario in which
projection-based methods and hypernetwork architectures have been most
thoroughly evaluated and in which structural comparisons between them are
most tractable.

% % ─────────────────────────────────────────────────────────────────────────────
\subsection{Evaluation Protocols: Multi-Head and Single-Head Designs}
\label{sec:eval_protocols}

Task-incremental benchmarks are evaluated under one of two output-layer
protocols. The \emph{multi-head} protocol gives each task a dedicated output
head, selected by task identity at test time, which removes output-level
interference between tasks; the \emph{single-head} protocol shares one output
layer, forcing the backbone alone to resolve that interference.
\citet{farquhar_towards_2018} showed that multi-head evaluation systematically
inflates apparent method performance by eliminating a major source of
forgetting, and argued that the single-head protocol is the more rigorous
test. Split-MNIST is therefore evaluated under both: the multi-head variant
reproduces the conditions of \citet{garg_fisher-orthogonal_2026} for direct
numerical comparison, and the single-head variant follows
\citet{farquhar_towards_2018} as the harder, more discriminating evaluation.
The exact head configuration and the protocol assigned to each benchmark are
given in Section~\ref{sec:benchmarks}.
% ─────────────────────────────────────────────────────────────────────────────
\subsection{Fisher Information Matrix}
\label{sec:fisher_rw}

Several methods surveyed in this thesis quantify parameter importance through
the \emph{Fisher Information Matrix} (FIM).
For a parametric model $p_\theta(y \mid x)$ with parameters
$\theta \in \mathbb{R}^P$, the FIM is the $P \times P$ matrix
\begin{equation}
    F(\theta)
    = \mathbb{E}_{x \sim q(x)}\,
      \mathbb{E}_{y \sim p_\theta(\cdot \mid x)}
      \Bigl[
        \nabla_\theta \log p_\theta(y \mid x)\,
        \nabla_\theta \log p_\theta(y \mid x)^\top
      \Bigr],
    \label{eq:fim}
\end{equation}
where $q(x)$ is the data-generating distribution \citep{amari_natural_1998}.
$F(\theta)$ defines a Riemannian metric on the parameter manifold: the
Fisher inner product $\langle u, v\rangle_F = u^\top F(\theta)\,v$ measures
distances in terms of the change in model predictions, not in terms of raw
Euclidean parameter displacement.

The full FIM is intractable at the scale of neural networks.
All methods evaluated in this thesis use the \emph{diagonal empirical
approximation}, which replaces the outer-product expectation with an
empirical average over $N$ observed data points:
\begin{equation}
    \hat{F}_i
    \;\approx\;
    \frac{1}{N}\sum_{n=1}^{N}
    \left(\frac{\partial \mathcal{L}^{(n)}}{\partial \theta_i}\right)^{\!2},
    \label{eq:empirical_fisher}
\end{equation}
yielding a non-negative scalar $\hat{F}_i$ for each parameter $\theta_i$.
The resulting vector $\hat{F} \in \mathbb{R}^P_{\geq 0}$ is treated as a
diagonal matrix in all subsequent quadratic forms.
This approximation reduces storage from $\mathcal{O}(P^2)$ to
$\mathcal{O}(P)$ and is the variant used by EWC
\citep{kirkpatrick_overcoming_2017} and all Fisher-based projection methods
evaluated in \citep{garg_fisher-orthogonal_2026}.

% ─────────────────────────────────────────────────────────────────────────────
\subsection{Regularization-Based Approaches}
\label{sec:rw_regularization}

The dominant paradigm for addressing catastrophic forgetting is
\emph{parameter regularization}, of which Elastic Weight Consolidation (EWC)
\citep{kirkpatrick_overcoming_2017} is the canonical representative.
EWC uses $F_i$ (Equation~\ref{eq:empirical_fisher}) as an importance
weight for each parameter, and applies a soft quadratic penalty that resists
changes to parameters deemed important to prior tasks:
\begin{equation}
\label{eq:ewc}
    \mathcal{L}_{\mathrm{EWC}}(\theta)
    = \mathcal{L}_t(\theta)
    + \frac{\lambda}{2}\sum_i F_i
      \bigl(\theta_i - \theta_{t-1,i}^*\bigr)^2,
\end{equation}
where $\theta \in \mathbb{R}^p$ are the model parameters,
$\theta^*_{t-1}$ is their snapshot saved after task $t{-}1$,
and $\lambda > 0$ governs overall regularization strength.
EWC is computationally efficient and broadly applicable, and serves as the
standard reference point against which new continual learning methods are
measured.
Its principal limitation is that soft penalties permit bounded parameter
drift: cumulative updates across a long task sequence gradually erode the
protection of earlier task spaces, a failure mode that becomes increasingly
pronounced as $T$ grows \citep{de_lange_continual_2022}.

% ─────────────────────────────────────────────────────────────────────────────
\subsection{Gradient Projection Methods}
\label{sec:rw_projection}

To enforce strict, non-negotiable protection for prior knowledge rather than
relying on penalties that can be overcome by sufficiently large gradients,
a second paradigm uses explicit \emph{gradient projection}.
Orthogonal Gradient Descent (OGD) \citep{farajtabar_orthogonal_2019} maintains
a memory matrix $G = [g_1 \mid \cdots \mid g_K] \in \mathbb{R}^{p \times K}$
of historical task gradient directions and projects incoming gradients onto the
orthogonal complement of $\operatorname{Col}(G)$ before each parameter update:
\begin{equation}
    \tilde{g} = g - G(G^\top G)^{-1}G^\top g.
\end{equation}
This provides a hard guarantee: the projected update cannot increase the loss
on any previously seen task.
Generalization bounds established within the Neural Tangent Kernel regime
confirm that OGD provides exact protection in linear and mildly nonlinear
settings \citep{bennani_generalisation_2020}.

However, standard Euclidean projection treats the parameter space as flat,
assigning equal importance to all directions regardless of the curvature of
the loss surface.

Fisher-Orthogonal Projected Natural Gradient Descent (FOPNG)
\citep{garg_fisher-orthogonal_2026} addresses this by performing the
projection in the Riemannian geometry defined by $\hat{F}$, the diagonal
empirical Fisher of the current task, so that the orthogonality constraint
respects the local curvature of the loss surface.
The related Fisher Natural Gradient (FNG) \citep{garg_fisher-orthogonal_2026}
applies the natural gradient without any projection, and Orthogonal Natural
Gradient (ONG) \citep{yadav_ong_2025} combines Fisher preconditioning with
Euclidean projection, sitting structurally between OGD and FOPNG.
Exact algorithmic instantiations are given in Section~\ref{sec:projection_methods}.

% ─────────────────────────────────────────────────────────────────────────────
\subsection{Hypernetworks for Continual Learning}
An orthogonal architectural response to catastrophic forgetting replaces the
multi-head output design entirely with a \emph{hypernetwork}
\citep{ha_hypernetworks_2017}: a meta-model
$\mathcal{H}(\cdot;\phi)$ trained to synthesize the complete weight
vector $\theta_t = \mathcal{H}(e_t;\phi)$ of a target network,
conditioned on a task embedding $e_t$.
Each task receives a unique weight configuration generated on demand, rather
than a dedicated output head, achieving implicit task separation through the
generative bottleneck.

Generating the full $\theta_t$ in a single forward pass is computationally
intractable for large target networks, so \citet{oswald_continual_2020}
introduced \emph{weight chunking}: the generator iterates over $C$ chunk
embeddings $c_1,\dots,c_C$, producing a fixed-size block of target parameters
per pass and concatenating them into the full weight vector $\theta_t$.
To prevent the hypernetwork from forgetting how to generate the weights of
earlier tasks, a soft \emph{functional regularizer} anchors the generator in
its own output space.
Let $\theta_s = \mathcal{H}(e_s;\phi)$ be the target-network weight vector the
generator currently produces for task $s$ under the single, evolving parameter
set $\phi$.
On completion of task $s$, the generated vector is cached as a fixed reference
$\theta^{\mathrm{ref}}_s$, detached from the computation graph.
While task $t$ is trained, only $\phi$ is updated, and every reference
$\{\theta^{\mathrm{ref}}_s\}_{s<t}$ stays fixed at the value it held when its
task finished, under the penalty
\begin{equation}
    \mathcal{L}_{\mathrm{reg}}
    = \frac{\beta}{t-1} \sum_{s=1}^{t-1}
      \bigl\| \mathcal{H}(e_s;\phi) - \theta^{\mathrm{ref}}_s \bigr\|_2^2,
\end{equation}
with regularization strength $\beta > 0$.
The penalty therefore measures, in the space of generated weights, how far the
current generator has drifted from the exact target parameters it once produced
for each past task.
This can be interpreted as EWC applied at the meta level: rather than anchoring individual target
parameters, it anchors the generated weight configurations for each previously
seen task, penalising drift in the function space where task-relevant
quantities actually live.
% ──────────────────
\subsection{Identified Gaps and Motivating Pathologies}

Whether hard geometric projection constraints are compatible with the
soft functional regularization already employed by hypernetworks, or
whether layering both simultaneously over-constrains the generator's
plasticity, has not been investigated (Sub-RQ1).

The architectural necessity of weight-chunking introduces a
gradient accumulation pathology --- inflated condition numbers in
the projection kernel --- that is unique to the hypernetwork and projection methods combination
and has no prior treatment in the literature (Sub-RQ2).

Existing Fisher-based projection methods define the projection geometry from a
single-task Fisher metric, which does not weight parameters by their importance
to past tasks; whether replacing it with a combined metric
$F_c = F_{\mathrm{new}} + F_{\mathrm{old}}$ that anchors historically important
coordinates, giving them inertia, yields consistent retention improvements
across benchmarks has not been examined (Sub-RQ3).

The historical Fisher estimate underlying Fisher-based projection
depends on a choice of accumulation operator; whether a monotonic
maximum or a decaying exponential moving average better preserves
long-horizon retention has not been directly compared (Sub-RQ4).